\section{Cronograma do Trabalho e Marcos do Projeto}

O desenvolvimento do Sistema de Gestão Hospitalar Inteligente Vitalis compreende um ciclo contínuo de 28 semanas (aproximadamente sete meses de execução técnica). A metodologia de trabalho organiza a construção do software de forma iterativa e incremental \cite{pressman:2016, sommerville:2019}, com as atividades divididas em seis fases estruturadas. Cada ciclo quinzenal de desenvolvimento entrega componentes funcionais avaliados por marcos de controle formais (\textit{milestones}), o que viabiliza a validação contínua dos requisitos clínicos e operacionais.

As seis fases do cronograma distribuem-se da seguinte forma:
\begin{enumerate}[leftmargin=*]
    \item \textbf{Fase 1: Concepção, elicitação e engenharia de requisitos} (Semanas 1 a 4);
    \item \textbf{Fase 2: Arquitetura de microsserviços e modelagem de dados 3NF} (Semanas 5 a 8);
    \item \textbf{Fase 3: Desenvolvimento do core clínico e triagem Manchester} (Semanas 9 a 14);
    \item \textbf{Fase 4: Módulo de inteligência artificial clínica e governança} (Semanas 15 a 20);
    \item \textbf{Fase 5: Gestão de leitos, faturamento TISS/TUSS e painéis gerenciais} (Semanas 21 a 24);
    \item \textbf{Fase 6: Testes integrados, simulações clínicas, segurança e homologação} (Semanas 25 a 28).
\end{enumerate}

\subsection{Estrutura Analítica das Fases e Marcos do Projeto}

Cada fase do projeto possui objetivos técnicos delimitados, atividades práticas e marcos de entrega que condicionam o avanço para a etapa seguinte.

\subsubsection{Fase 1: Concepção, elicitação e engenharia de requisitos (Semanas 1 a 4)}

A etapa inicial concentra o alinhamento com os usuários do hospital e a formalização das especificações técnicas:
\begin{itemize}[leftmargin=*]
    \item \textbf{Semanas 1 e 2}: Condução de entrevistas estruturadas com médicos emergencistas, enfermeiros triadores, gestores de hotelaria hospitalar e faturistas. Consolidação dos 22 requisitos funcionais, 12 regras de negócio e requisitos não funcionais de desempenho e segurança na Especificação de Requisitos de Software (ERS).
    \item \textbf{Semanas 3 e 4}: Mapeamento completo dos processos de trabalho hospitalares em diagramas BPMN 2.0 (acolhimento, classificação de risco, internação, prescrição eletrônica e faturamento). Análise de conformidade legal com a Lei Geral de Proteção de Dados Pessoais (LGPD) e com as resoluções do Conselho Federal de Medicina (CFM nº 1.821/2007 e nº 2.299/2021).
\end{itemize}

Os marcos de controle definidos para a Fase 1 são:
\begin{itemize}[leftmargin=*]
    \item \textbf{Marco M1.1 (Semana 2)}: Especificação de Requisitos de Software (ERS) homologada com matriz preliminar de rastreabilidade.
    \item \textbf{Marco M1.2 (Semana 4)}: Modelagem de processos em BPMN 2.0 e relatório de diretrizes de conformidade com a LGPD concluídos.
\end{itemize}

\subsubsection{Fase 2: Arquitetura de microsserviços e modelagem de dados 3NF (Semanas 5 a 8)}

A segunda etapa estabelece os alicerces estruturais de software e persistência de dados:
\begin{itemize}[leftmargin=*]
    \item \textbf{Semanas 5 e 6}: Modelagem e implementação do banco de dados relacional PostgreSQL em Terceira Forma Normal (3NF). O esquema abrange entidades centrais (Paciente, Atendimento, Prontuário Lógico, Evolução, Prescrição, Leito e Faturamento), com integridade referencial estrita e tabelas de log imutáveis para auditoria clínica e inferências de inteligência artificial.
    \item \textbf{Semanas 7 e 8}: Estruturação do barramento de microsserviços com NestJS e Fastify, conteinerização dos módulos via Docker e orquestração básica. Configuração de filas assíncronas com Redis e BullMQ para processamento de eventos hospitalares, esteira de integração contínua (CI/CD) e módulo de controle de acesso híbrido (RBAC e ABAC).
\end{itemize}

Os marcos de controle definidos para a Fase 2 são:
\begin{itemize}[leftmargin=*]
    \item \textbf{Marco M2.1 (Semana 6)}: Esquema relacional 3NF implementado com scripts de migração automatizados e integridade referencial testada.
    \item \textbf{Marco M2.2 (Semana 8)}: Infraestrutura de microsserviços, mensageria distribuída e esteira de CI/CD operacionais no ambiente de desenvolvimento.
\end{itemize}

\subsubsection{Fase 3: Desenvolvimento do core clínico e triagem Manchester (Semanas 9 a 14)}

A terceira etapa foca na construção das ferramentas de atendimento direto ao paciente:
\begin{itemize}[leftmargin=*]
    \item \textbf{Semanas 9 a 11}: Construção do módulo de recepção com validação cadastral de CPF e integração ao Cartão Nacional de Saúde (CNS). Desenvolvimento da interface de triagem com implementação determinística do Protocolo de Manchester (cinco níveis de prioridade), formulário de coleta de sinais vitais e transcrição de áudio via API para queixas dos pacientes.
    \item \textbf{Semanas 12 a 14}: Desenvolvimento do Prontuário Eletrônico do Paciente (PEP), com garantia de imutabilidade dos registros após a assinatura médica. Implementação da prescrição médica digital estruturada, incluindo checagem obrigatória de alergias cadastradas e suporte à assinatura digital padrão ICP-Brasil (PAdES e PKCS\#7).
\end{itemize}

Os marcos de controle definidos para a Fase 3 são:
\begin{itemize}[leftmargin=*]
    \item \textbf{Marco M3.1 (Semana 11)}: Módulo de recepção e triagem Manchester concluído com reorganização automática de filas para casos classificados como Emergência (Vermelho).
    \item \textbf{Marco M3.2 (Semana 14)}: Prontuário eletrônico imutável e prescrição médica com verificação de alergias e assinatura digital em pleno funcionamento.
\end{itemize}

\subsubsection{Fase 4: Módulo de inteligência artificial clínica e governança (Semanas 15 a 20)}

A quarta etapa integra as capacidades de inteligência artificial preditiva e generativa sob salvaguardas éticas:
\begin{itemize}[leftmargin=*]
    \item \textbf{Semanas 15 a 17}: Construção do gateway de anonimização e sanitização de dados pessoais identificáveis (PII), com aplicação de tokenização e $k$-anonimato ($k \ge 5$) antes de qualquer chamada a APIs externas. Desenvolvimento do motor preditivo para cálculo contínuo dos escores de instabilidade hemodinâmica e sepse precoce (qSOFA e NEWS2) vinculado aos dados da triagem e internação.
    \item \textbf{Semanas 18 a 20}: Construção do copiloto clínico de prescrição com arquitetura RAG (\textit{Retrieval-Augmented Generation}), utilizando banco vetorial alimentado com literatura médica, protocolos institucionais e bulários \cite{lewis_rag:2020}. Implementação da síntese automática de evoluções médicas e minutas de sumários de alta hospitalar via modelos de linguagem, sob supervisão médica estrita (\textit{Human-in-the-Loop}). Configuração do mecanismo de \textit{Circuit Breaker} para acionamento de contingência local quando a latência de rede superar 2,0 segundos.
\end{itemize}

Os marcos de controle definidos para a Fase 4 são:
\begin{itemize}[leftmargin=*]
    \item \textbf{Marco M4.1 (Semana 17)}: Gateway de sanitização PII aprovado em testes de vazamento de dados e motor preditivo qSOFA/NEWS2 respondendo em menos de 1,5 segundos.
    \item \textbf{Marco M4.2 (Semana 20)}: Copiloto farmacoterapêutico baseado em RAG e gerador de sumários clínicos homologados com interface de revisão médica.
\end{itemize}

\subsubsection{Fase 5: Gestão de leitos, faturamento TISS/TUSS e painéis gerenciais (Semanas 21 a 24)}

A quinta etapa contempla as operações logísticas, administrativas e financeiras da instituição:
\begin{itemize}[leftmargin=*]
    \item \textbf{Semanas 21 e 22}: Desenvolvimento do mapa interativo de leitos com atualização reativa via WebSockets e controle dos estados de Livre, Ocupado, Higienização, Manutenção e Isolamento. Acoplamento do modelo de aprendizado de máquina para estimativa do tempo provável de permanência (\textit{Length of Stay} ou LOS) e apoio à regulação de vagas.
    \item \textbf{Semanas 23 e 24}: Construção do módulo de faturamento hospitalar para geração de lotes eletrônicos no padrão TISS/TUSS da ANS \cite{ans_tiss:2023}. Implementação da auditoria preventiva por inteligência artificial para detecção precoce de inconsistências documentais e glosas financeiras. Desenvolvimento de painéis de telemetria hospitalar (tempo médio de espera, taxa de ocupação e giro de leitos) com exportação em PDF, CSV e JSON FHIR \cite{hl7_fhir:2019}.
\end{itemize}

Os marcos de controle definidos para a Fase 5 são:
\begin{itemize}[leftmargin=*]
    \item \textbf{Marco M5.1 (Semana 22)}: Painel de leitos sincronizado em tempo real e modelo preditivo de tempo de permanência ativos.
    \item \textbf{Marco M5.2 (Semana 24)}: Módulo de faturamento TISS com auditoria preventiva de glosas e painéis gerenciais operando de forma integrada.
\end{itemize}

\subsubsection{Fase 6: Testes integrados, simulações clínicas, segurança e homologação (Semanas 25 a 28)}

A etapa final submete a plataforma hospitalar a verificações de confiabilidade sob estresse:
\begin{itemize}[leftmargin=*]
    \item \textbf{Semanas 25 e 26}: Bateria de testes de carga com 2.000 conexões concorrentes para certificar a estabilidade transacional do banco de dados e a resposta das APIs. Realização de testes de penetração (\textit{pentests}), auditoria de segurança da informação (criptografia AES-256 e TLS 1.3) e validação da imutabilidade das trilhas de auditoria.
    \item \textbf{Semanas 27 e 28}: Execução de simulações clínicas em ambiente hospitalar de testes, com reprodução de fluxos de emergência médica, transferências de pacientes e falhas deliberadas de conectividade para testar o chaveamento determinístico. Capacitação técnica de multiplicadores (médicos, enfermeiros e administradores) e liberação formal do Sistema Vitalis.
\end{itemize}

Os marcos de controle definidos para a Fase 6 são:
\begin{itemize}[leftmargin=*]
    \item \textbf{Marco M6.1 (Semana 26)}: Testes de estresse concluídos com atendimento aos índices de latência e relatório de auditoria de segurança sem vulnerabilidades críticas.
    \item \textbf{Marco M6.2 (Semana 28)}: Homologação clínica completa realizada, equipe treinada e entrega formal do sistema para operação piloto.
\end{itemize}

\subsection{Matriz de Atividades, Prazos e Entregáveis do Projeto}

A Tabela \ref{tab:cronograma_vitalis} sintetiza a distribuição das 28 semanas do projeto, relacionando as fases de engenharia aos seus componentes, semanas de execução e entregáveis formais.

\begin{table}[H]
\centering
\scriptsize
\caption{Estrutura analítica de atividades, cronograma e entregáveis do Vitalis}
\label{tab:cronograma_vitalis}
\begin{tabularx}{\textwidth}{l >{\raggedright\arraybackslash}p{2.8cm} >{\raggedright\arraybackslash}X c >{\raggedright\arraybackslash}p{3.5cm}}
\toprule
\textbf{Fase} & \textbf{Componente} & \textbf{Atividades Principais} & \textbf{Semanas} & \textbf{Entregável Formal} \\
\midrule
\multirow{2}{*}{\textbf{Fase 1}} & Elicitação e Requisitos & Entrevistas com especialistas, levantamento de escopo e redação da ERS. & 1 a 2 & Documento formal de ERS  aprovado. \\
\cmidrule(lr){2-5}
 & Processos e Regulação & Modelagem de processos em BPMN 2.0 e análise de conformidade com a LGPD e CFM. & 3 a 4 & Diagramas BPMN e relatório de conformidade jurídica. \\
\midrule
\multirow{2}{*}{\textbf{Fase 2}} & Modelagem de Dados & Normalização relacional em 3NF e criação de tabelas de log auditáveis. & 5 a 6 & DDL PostgreSQL com migrações versionadas. \\
\cmidrule(lr){2-5}
 & Microsserviços e Filas & Barramento NestJS, orquestração Docker, filas Redis/BullMQ e autenticação RBAC/ABAC. & 7 a 8 & Repositório base com esteira CI/CD e mensageria funcional. \\
\midrule
\multirow{2}{*}{\textbf{Fase 3}} & Recepção e Manchester & Cadastro integrado ao CNS e motor determinístico da triagem de Manchester. & 9 a 11 & Módulo de acolhimento e fila prioritária reativa. \\
\cmidrule(lr){2-5}
 & Prontuário e Prescrição & PEP estruturado com imutabilidade e prescrição médica com assinatura ICP-Brasil. & 12 a 14 & Módulo PEP integrado e prescrição médica digital assinada. \\
\midrule
\multirow{2}{*}{\textbf{Fase 4}} & Sanitização e Risco & Gateway de sanitização PII e motor preditivo de instabilidade clínica (qSOFA/NEWS2). & 15 a 17 & Pipeline de desidentificação e alertas de sepse em tempo real. \\
\cmidrule(lr){2-5}
 & Copiloto RAG e LLM & Base vetorial clínica, recomendação posológica e síntese supervisionada de sumários. & 18 a 20 & Copiloto farmacológico e gerador de sumários homologados. \\
\midrule
\multirow{2}{*}{\textbf{Fase 5}} & Leitos e Predição LOS & Painel visual de leitos via WebSockets e predição do tempo de permanência hospitalar. & 21 a 22 & Mapa de leitos interativo e modelo preditivo de alta. \\
\cmidrule(lr){2-5}
 & Faturamento e Painéis & Geração de lotes XML TISS/TUSS, auditoria preditiva de glosas e painéis gerenciais. & 23 a 24 & Módulo de contas médicas e painel de indicadores operacionais. \\
\midrule
\multirow{2}{*}{\textbf{Fase 6}} & Carga e Segurança & Testes com 2.000 requisições simultâneas, pentests e validação de criptografia. & 25 a 26 & Relatório de conformidade de carga e segurança da informação. \\
\cmidrule(lr){2-5}
 & Simulações e Release & Ensaios clínicos de emergência, treinamento de usuários e homologação do sistema. & 27 a 28 & Pacote de implantação homologado e manuais técnicos. \\
\bottomrule
\end{tabularx}
\end{table}

\subsection{Representação Visual do Cronograma (Gráfico de Gantt)}

A Figura \ref{fig:gantt_vitalis} ilustra a linha temporal das 28 semanas do projeto por meio de um gráfico de Gantt elaborado com o pacote \texttt{pgfgantt}. A visualização destaca o sequenciamento das tarefas e os marcos críticos de encerramento de cada etapa.

\begin{figure}[H]
\centering
\resizebox{\textwidth}{!}{
\begin{ganttchart}[
    canvas/.append style={fill=none, draw=vitalisdark!30, line width=.75pt},
    hgrid style/.style={draw=vitalisdark!12, line width=.4pt},
    vgrid={*1{draw=vitalisdark!12, line width=.4pt}},
    x unit=0.46cm,
    y unit chart=0.68cm,
    y unit title=0.65cm,
    title/.style={draw=vitalisdark!30, fill=vitalislight!70},
    title label font=\bfseries\scriptsize\color{vitalisdark},
    bar/.append style={fill=vitalisblue!80, draw=vitalisdark, line width=0.5pt, rounded corners=2pt},
    bar label font=\scriptsize\sffamily,
    bar height=0.52,
    group/.append style={fill=vitalisdark!85, draw=vitalisdark},
    group label font=\bfseries\scriptsize\sffamily\color{vitalisdark},
    group height=0.22,
    milestone/.append style={fill=vitalisred, draw=vitalisred!80!black},
    milestone label font=\scriptsize\bfseries\sffamily\color{vitalisred!80!black},
    milestone height=0.4
]{1}{28}
    % Títulos das colunas
    \gantttitle{Cronograma Global de Desenvolvimento: Sistema Vitalis (28 Semanas)}{28} \\
    \gantttitle{Fase 1}{4}
    \gantttitle{Fase 2}{4}
    \gantttitle{Fase 3}{6}
    \gantttitle{Fase 4}{6}
    \gantttitle{Fase 5}{4}
    \gantttitle{Fase 6}{4} \\
    \gantttitlelist{1,...,28}{1} \\

    % Fase 1
    \ganttgroup{Fase 1: Requisitos e Processos}{1}{4} \\
    \ganttbar{Elicita\c{c}\~ao de Requisitos e ERS}{1}{2} \\
    \ganttbar{Mapeamento BPMN e Parecer LGPD}{3}{4} \\
    \ganttmilestone{M1: Requisitos Homologados}{4} \\

    % Fase 2
    \ganttgroup{Fase 2: Arquitetura e Modelagem 3NF}{5}{8} \\
    \ganttbar{Banco de Dados Relacional 3NF}{5}{6} \\
    \ganttbar{Microsservi\c{c}os, Redis e CI/CD}{7}{8} \\
    \ganttmilestone{M2: Infraestrutura Operacional}{8} \\

    % Fase 3
    \ganttgroup{Fase 3: Core Cl\'inico e Manchester}{9}{14} \\
    \ganttbar{Recep\c{c}\~ao e Triagem Manchester}{9}{11} \\
    \ganttbar{PEP e Prescri\c{c}\~ao com Assinatura}{12}{14} \\
    \ganttmilestone{M3: Fluxo Assistencial Ativo}{14} \\

    % Fase 4
    \ganttgroup{Fase 4: M\'odulo de IA e Governan\c{c}a}{15}{20} \\
    \ganttbar{Sanitizador PII e Predi\c{c}\~ao qSOFA}{15}{17} \\
    \ganttbar{Copiloto RAG e S\'intese Cl\'inica}{18}{20} \\
    \ganttmilestone{M4: Pipeline de IA Homologado}{20} \\

    % Fase 5
    \ganttgroup{Fase 5: Leitos e Faturamento TISS}{21}{24} \\
    \ganttbar{Mapa de Leitos e Estimativa LOS}{21}{22} \\
    \ganttbar{Faturamento TISS/TUSS e Pain\'eis}{23}{24} \\
    \ganttmilestone{M5: Gest\~ao Hospitalar Pronta}{24} \\

    % Fase 6
    \ganttgroup{Fase 6: Testes, Seguran\c{c}a e Release}{25}{28} \\
    \ganttbar{Testes de Carga e Seguran\c{c}a (Pentest)}{25}{26} \\
    \ganttbar{Simula\c{c}\~oes Cl\'inicas e Capacita\c{c}\~ao}{27}{28} \\
    \ganttmilestone{M6: Release Final Vitalis}{28}
\end{ganttchart}
}
\caption{Gráfico de Gantt representativo das 28 semanas do projeto Vitalis}
\label{fig:gantt_vitalis}
\end{figure}

\subsection{Caminho Crítico e Gestão de Riscos do Projeto}

O sucesso do cronograma depende do monitoramento contínuo das dependências técnicas que compõem o caminho crítico (\textit{Critical Path Method}). O encadeamento estrito do projeto exige que cada fase forneça os artefatos fundamentais para a subsequente:
\begin{enumerate}[leftmargin=*]
    \item A conclusão da modelagem de dados relacional 3NF (Fase 2) depende diretamente do refinamento dos atributos contidos na ERS (Fase 1);
    \item A gravação de prontuários eletrônicos e prescrições médicas (Fase 3) exige a disponibilidade das tabelas relacionais com integridade referencial ativa (Fase 2);
    \item A integração dos motores de inteligência artificial (Fase 4) depende da maturidade dos dados clínicos gerados na triagem e no prontuário (Fase 3), além da operação confiável do gateway de sanitização de dados pessoais identificáveis (PII);
    \item A auditoria preventiva de contas médicas (Fase 5) requer o cruzamento entre as prescrições médicas executadas (Fase 3) e os modelos de detecção de inconformidades (Fase 4);
    \item A bateria final de testes de carga e simulação clínica (Fase 6) pressupõe todos os módulos integrados no mesmo ambiente operacional.
\end{enumerate}

Para mitigar desvios durante as 28 semanas, foram estabelecidas quatro contingências preventivas:
\begin{itemize}[leftmargin=*]
    \item \textbf{Instabilidade ou latência de APIs externas de IA}: O sistema conta com um componente de \textit{Circuit Breaker}. Se a latência de rede exceder 2,0 segundos ou houver indisponibilidade momentânea da nuvem de IA, o sistema chaveia automaticamente para a avaliação determinística local (tabelas convencionais de Manchester e cálculo estático de doses), sem bloquear o atendimento médico.
    \item \textbf{Vazamento inadvertido de dados sensíveis}: O pipeline de sanitização executa testes de regressão automatizados em todas as alterações de código para garantir que nenhuma entidade identificável (CPF, nome, registro geral ou endereço) seja transmitida externamente.
    \item \textbf{Rejeição de arquivos de faturamento pelas operadoras}: A validação sintática contra os esquemas XSD oficiais da ANS é realizada no próprio ambiente hospitalar antes do fechamento do lote XML, o que previne retrabalhos com glosas administrativas.
    \item \textbf{Resistência profissional ao uso de recomendações automatizadas}: A interface adota o princípio de supervisão médica mandatória (\textit{Human-in-the-Loop}). Nenhum escore de risco ou sugestão farmacológica entra em vigor sem a aprovação explícita e assinatura digital do médico ou enfermeiro responsável.
\end{itemize}

\destaqueVitalis{Marco Final e Entrada em Operação Piloto (Semana 28)}{%
A conclusão bem-sucedida da Semana 28 culmina na entrega formal e homologação do Sistema Vitalis. O sistema entra em operação assistida com a equipe multiprofissional treinada, documentação técnica homologada, cobertura de testes de carga com 2.000 conexões concorrentes e conformidade integral com a LGPD e as resoluções do Conselho Federal de Medicina.
}
